\documentclass[11pt, a4paper]{article}

\usepackage[a4paper, top=2.5cm, bottom=2.5cm, left=2cm, right=2cm]{geometry}
\usepackage{amsmath}
\usepackage{amssymb}
\usepackage[colorlinks=true, linkcolor=blue, urlcolor=blue, citecolor=blue]{hyperref}

\title{Theory of Everything - Volume 43 \\ \Large Galactic Ecosystems}
\author{Omega Protocol}
\date{\today}

\begin{document}

\maketitle

\section*{Layman's Summary}
What does a civilization that controls a whole galaxy look like? Volume 43 proves that they wouldn't use spaceships or laser beams. A Type III civilization would manipulate the very fabric of spacetime, using supermassive black holes as their primary computers and energy sources.

\section{Galactic Network Hubs (Axiom 46)}
Every large galaxy has a supermassive black hole at its center. We establish that for a Type III civilization, this black hole is not a hazard, but the ultimate prize. The supermassive black hole serves as the central "Archive node"—a gargantuan, perfectly efficient quantum computer anchoring the civilization's network.

\section{Penrose Energy Extraction (Theorem)}
How do you power a galactic computer? We prove that civilizations can use the "Penrose Process" to extract near-infinite computational energy from the "ergosphere"—the swirling tornado of spacetime just outside a spinning black hole. You drop garbage in, and extract pure, free energy out.

\section{The Interstellar Internet (Theorem)}
Radio waves take thousands of years to cross a galaxy; they are useless for a unified civilization. We prove that a Type III civilization must rely exclusively on "entangled photon streams." By exploiting the "Archive Mode" of the universe, they can transfer information instantly across galactic distances without violating relativity.

\end{document}
